\documentclass[12pt]{article}

\usepackage{setspace}
\onehalfspace

%\usepackage[utf8]{inputenc}
%\usepackage{t1enc}
%\usepackage[T1]{fontenc}

\usepackage{amsmath}

\usepackage{moodle}

\DeclareMathOperator{\Exp}{Exp}
\DeclareMathOperator{\Norm}{{\cal N}}
\DeclareMathOperator{\Binom}{Binom}
\DeclareMathOperator{\E}{E}
\DeclareMathOperator{\D}{D}
\renewcommand{\P}{\operatorname{P}}
%\DeclareMathOperator{\cov}{cov}
\DeclareMathOperator{\corr}{corr}

\DeclareMathOperator{\F}{F} 
\renewcommand{\d}{\operatorname{d\! }} %integral \! a nyero


\DeclareMathOperator{\R}{\mathbb{R}}
%\DeclareMathOperator{\Norm}{{\cal N}}
\DeclareMathOperator{\Unif}{{\cal U}}
\DeclareMathOperator{\Poi}{{Poisson}}


\DeclareMathOperator{\cov}{cov}





\begin{document}

\begin{quiz}{gyakorló}

\begin{multi}{esem10}
Ha $\P(A)=0$ akkor $A$ sosem következik be.
\item igaz
\item* hamis
\end{multi}
\begin{multi}{impdvv2}
Egy $\xi\sim \Binom(n,p)$ várható értéke: $n+p$
\item igaz
\item* hamis
\end{multi}
\begin{multi}{impfvv4}
Egy $\xi\sim \Unif(a,b)$ szórásnégyzete: $\frac{(a+b)^2}{2}$
\item igaz
\item* hamis
\end{multi}
\begin{multi}{flenvv7}
Ha $\xi$ és $\eta$ valségi változóknak létezik a szórása, akkor a 
kovarianciájuk
$$
\cov(\xi,\eta)=\E((\xi-\E(\xi))(\eta-\E(\eta)))
$$
\item* igaz
\item hamis
\end{multi}
\begin{multi}{felt4}
$A$-nak a $B$-re vonatkozó feltételes valségét $\P(A|B)=\frac{\P(A\cdot B)}{\P(B)}$ módon értelmezzük, feltéve hogy $\P(B)>0$.
\item* igaz
\item hamis
\end{multi}
\begin{multi}{eofv3}
Egy $F:\R\to [0,1]$ pontosan akkor eloszlásfüggvény, ha monoton nemcsökkenő, 
$\lim_{x\to+\infty}F(x)=1$ és $\lim_{x\to-\infty}F(x)=0$
\item igaz
\item* hamis
\end{multi}
\begin{multi}{sfv9}
Ha $\xi$-nek $f$ a sűrűségfüggvénye és $F$ az eloszlásfüggvénye, akkor bármely $b$-re:
$$
F(b)=\int_{\infty}^{b}f(x)\d x
$$
\item* igaz
\item hamis
\end{multi}
\begin{multi}{val1}
A valószínűség additív, azaz $\P(A+B)=\P(A)+\P(B)$.
\item igaz
\item* hamis
\end{multi}
\begin{multi}{felt3}
$A$-nak a $B$-re vonatkozó feltételes valsége $\P(A|B)=\frac{\P(A\cdot B)}{\P(B)}$.
\item igaz
\item* hamis
\end{multi}
\begin{multi}{eofv7}
Bármely $\xi$-re és $a<b$-re $\P(a\le \xi <b)=F_{\xi}(b)-F_{\xi}(a)$.
\item* igaz
\item hamis
\end{multi}
\begin{multi}{dvv2}
Egy nemnegatív tagú $p_1,\ldots, p_n$ vektor pont akkor valségeloszlás, ha $\sum_{k=1}^n p_k=1$.
\item* igaz
\item hamis
\end{multi}
\begin{multi}{dvv16}
$\xi$ és $\eta$ véges értékkészletű valségi változókra:
$$
\D^2(\xi+\eta)=\D^2(\xi)+\D^2(\eta)
$$
\item igaz
\item* hamis
\end{multi}
\begin{multi}{flenvv9}
Ha a $\xi$ és $\eta$ valségi változóknak létezik a korrelációja, akkor $-1\le \corr(\xi,\eta)\le 1$.
\item* igaz
\item hamis
\end{multi}
\begin{multi}{val5}
Bármely esemény valsége kiszámolható a $\frac{\text{kedvező}}{\text{összes}}$ képlettel.
\item igaz
\item* hamis
\end{multi}
\begin{multi}{felt8}
Ha $A_1,\ldots,A_n$ teljes eseményrendszer és $\P(B)>0$ akkor:
$$
\P(B)=\sum_{k=1}^n \P(A_k|B)\P(B)
$$
\item* igaz
\item hamis
\end{multi}
\begin{multi}{impfvv6}
Ha $\xi\sim \Unif(0,1)$ akkor  minden $a,b$-re $b\xi+a\sim \Unif(a,b)$.
\item igaz
\item* hamis
\end{multi}
\begin{multi}{sfv8}
Ha van $\xi$-nek sűrűségfüggvénye, akkor bármely $b$-re:
$$
\P(\xi<b)=f(b)
$$
\item igaz
\item* hamis
\end{multi}
\begin{multi}{impfvv10}
Ha $\xi\sim \Exp(\lambda)$ akkor a sűrűségfüggvénye 
$$
f_{\xi}(x)=
\begin{cases}
1-e^{-\lambda x}& \ \ x>0\\
0 & \text{máskor} 
\end{cases}
$$
\item igaz
\item* hamis
\end{multi}
\begin{multi}{sfv7}
Ha $\xi$-nek $f$ asűrűségfüggvénye, akkor bármely $a<b$-re:
$$
\P(a<\xi<b)=f(b)-f(a)
$$
\item igaz
\item* hamis
\end{multi}
\begin{multi}{esem8}
Az $A$ és $B$ események pont akkor függetlenek, ha $\P(A+B)=\P(A)+\P(B)$.
\item igaz
\item* hamis
\end{multi}
\begin{multi}{esem4}
Az $\overline{A} + \overline{B}$ maga után vonja az $\overline{A\cdot B}$ eseményt.
\item* igaz
\item hamis
\end{multi}
\begin{multi}{impdvv4}
Egy $\xi\sim \Binom(n,p)$ szórásnégyzete: $\frac{n+1}{2}$
\item igaz
\item* hamis
\end{multi}
\begin{multi}{dvv8}
Egy $\xi$ valségi változó az $x_1,\ldots x_n$ értékeket veheti fel és $p_k=\P(\xi=x_k)$. 
Ekkor a második momentuma:
$$
\E(\xi^2)=\sum_{k=1}^n x_k^2 p_k
$$
\item* igaz
\item hamis
\end{multi}
\begin{multi}{dvv6}
Egy $\xi$ valségi változó az $x_1,\ldots x_n$ értékeket veheti fel. Ekkor a szórása:
$$
\D(\xi)=\sum_{k=1}^n \P(\xi=x_k)(x_k-\E(\xi))^2
$$
\item igaz
\item* hamis
\end{multi}
\begin{multi}{eofv11}
Bármely $\xi$-re és $a$-ra $\P(\xi= a)=0$, feltéve hogy $F_{\xi}$ folytonos.
\item* igaz
\item hamis
\end{multi}
\end{quiz}

\end{document}

